\section{课内知识联动分析}
本次实习中涉及了大量的课内知识，本节将结合现场观察进行专题联动分析。\par{}

\subsection{上海地铁21号线车下设备分析}
上海地铁21号线采用$\mathrm{{T}_{c1}-{M}_{p1}-{M}_1-{M}_2-{M}_{p2}-{T}_{c2}}$的4动2拖布局，其$\mathrm{M}$车车下设备布局如图\ref{fig:l21-m-underfloor-layout}所示。\par{}
\reportFigure{17}
由图可得，该车从一位端到二位端依次配备了接线盒、转向架（未安装）、空气制动设备、牵引逆变器、低压箱、辅助逆变器、制动电阻与辅助制动设备、空气制动设备、转向架（未安装）、接线盒。\par{}
该车同时配备制动电阻与制动缸，说明该车采用电阻制动加空气制动的形式进行制动；根据牵引传动课程相关知识，由于电阻制动存在低速时电机反电势降低、可提供的电制动力逐渐减小的特点，仅依靠电阻制动无法制停列车，故需要通过空气制动进行低速制动；一般现代地铁的空气制动采用电信号进行控制，浦镇造列车大都使用盘形制动，其中又分为轮盘式和轴盘式，地铁多用轮盘式。\par{}
与21号线的$\mathrm{M_p}$车相比，虽然其都有牵引逆变器等核心牵引设备，但由于$\mathrm{M}$车不带受电弓，其不存在$\mathrm{M_p}$车的三位开关、断路器等设备；与$\mathrm{T_c}$车相比，$\mathrm{M}$车存在牵引逆变器但没有车头驾驶室结构。\par{}
同校内1动1拖编组的试验车相比，该车厢$\mathrm{M}$车自带一个辅助逆变器和低压箱，而学校里的试验车的辅逆安装在$\mathrm{T_c}$车上。值得注意的是靠近二位端的电空制动装置旁边存在着一片未安装设备的区域，尚不能确定该区域是否预留给后续安装的设备。\par{}
车下某些电气设备上写有如下几行字：\par{}
\reportFigure{18}
\noindent{}这些要求的目的是保证维修人员的安全，因此要求断电、三位开关接地、电荷通过放电电阻释放完毕等。\par{}
此外，由铭牌可得$\mathrm{M}$车上的辅助逆变器、低压箱、牵引逆变器，以及$\mathrm{M_p}$车上的高压箱等均由上海阿尔斯通公司生产，可见该公司业务的广泛。\par{}

\subsection{车体轻量化策略分析}
在轨道车辆设计课程中我们已经了解：车体轻量化存在两种主要方式，一是采用新型材料，二是采用先进设计方法。\par{}
材料方面，可以从材料成分、晶体形态等角度入手。如采用耐候钢等高耐久材料，以及铝合金、镁合金甚至碳纤维等轻量化材料；也可以从材料的宏观结构，如中空结构、复合材料等方向入手。讲座中，我们了解到浦镇公司正积极探索碳纤维在车体结构上的应用，深入研究其在结构、疲劳、应力分布等方面的作用机理，提升新型轻量化材料在车辆设计上的应用比例。\par{}
在设计方法方面，目前数字化设计与优化分析手段已十分成熟。通过广泛使用计算机辅助技术（CAx）、面向产品生命周期的技术（DFx）建立虚拟样机，针对数字化模型进行静力学仿真、动力学仿真、振动模态分析、失稳模态分析、疲劳分析等，可以节省制造实体模型并进行测试的成本；通过尺寸优化、形貌优化、拓扑优化等方法，在某些约束下可以使得某个或某些设计变量达到极值，从而实现减重的目的。通过浦镇公司的讲座还了解到，当下人工智能技术发展迅速，可以采用AI+自动化设计与仿真的手段，在给定边界、应力、重量和模态多目标优化条件下，让智能系统给出多种优化结构，供设计师进行初选和思路发散，提升设计新颖性与实用性。\par{}

